\section{A brief review of Hodge modules}\label{Hodge_modules}
A key ingredient for the definition of our invariants is Saito's theory of mixed Hodge modules. In what follows,
we give a brief presentation of the relevant objects, and recall a few facts that we will need.
For details, we refer to \cite{Saito-MHM}.

Given a smooth $n$-dimensional complex algebraic variety $X$, we denote by $\jepPubScript{D}_X$ the sheaf of differential operators on $X$. 
This carries the increasing filtration $F_{\jepPubBullet}\jepPubScript{D}_X$ by order of differential operators. A \emph{left} or \emph{right $\jepPubScript{D}$-module}
is a left, respectively right, $\jepPubScript{D}_X$-module, which is quasi-coherent as an $\jepPubScript{O}_X$-module. There is an equivalence 
between the categories of left and right $\jepPubScript{D}$-modules, which at the level of $\jepPubScript{O}_X$-modules is given by
$$\jepPubScript{M}\longrightarrow \jepPubScript{N}:=\jepPubScript{M}\otimes_{\jepPubScript{O}_X}\omega_X\quad\text{and}\quad \jepPubScript{N}\longrightarrow \jepPubScript{H}om_{\jepPubScript{O}_X}(\omega_X,\jepPubScript{N}).$$
For example, this equivalence maps the left $\jepPubScript{D}$-module $\jepPubScript{O}_X$ to the right $\jepPubScript{D}$-module $\omega_X$.
For a thorough introduction to the theory of $\jepPubScript{D}$-modules, we refer to \cite{HTT}.

A \emph{filtered left} (or \emph{right}) \emph{$\jepPubScript{D}$-module} 
is a $\jepPubScript{D}$-module $\jepPubScript{M}$, together with an increasing filtration $F=F_{\jepPubBullet}\jepPubScript{M}$ that is compatible with the
order filtration on $\jepPubScript{D}_X$ and which is \emph{good}, in a sense to be defined momentarily. A morphism of filtered $\jepPubScript{D}$-modules is required to be compatible with the filtrations. The equivalence between left
and right $\jepPubScript{D}$-modules extends to the categories of filtered modules, with the convention that 
$$F_{p-n}(\jepPubScript{M}\otimes_{\jepPubScript{O}_X}\omega_X)=F_p\jepPubScript{M}\otimes_{\jepPubScript{O}_X}\omega_X.$$
A filtration $F_{\jepPubBullet}\jepPubScript{M}$ on a coherent $\jepPubScript{D}$-module $\jepPubScript{M}$ is \emph{good} if the corresponding graded object 
${\rm gr}_{\jepPubBullet}^F\jepPubScript{M}:=\mathop{\textstyle\bigoplus}\limits_kF_k\jepPubScript{M}/F_{k-1}\jepPubScript{M}$ is locally finitely generated over ${\rm gr}_{\jepPubBullet}^F\jepPubScript{D}_X$. 
We note that
every coherent $\jepPubScript{D}$-module admits a good filtration, but this is far from being unique. 


We now come to the key objects in Saito's theory, the \emph{mixed Hodge modules} from \cite{Saito-MHM}. Such an object is given by the data 
$M=(\jepPubScript{M},F,\jepPubScript{P} ,\varphi, W)$, where:
\begin{enumerate}
\item[(i)] $(\jepPubScript{M},F)$ is a filtered $\jepPubScript{D}$-module, with $\jepPubScript{M}$ a holonomic left (or right) $\jepPubScript{D}$-module, with regular singularities; 
$F$ is the \emph{Hodge filtration} of $\jepPubScript{M}$.
\item[(ii)] $\jepPubScript{P}$ is a perverse sheaf of $\jepPubBold{Q}$-vector spaces on $X$.
\item[(iii)] $\varphi$ is an isomorphism between $\jepPubScript{P}_{\jepPubBold{C}} = \jepPubScript{P} \otimes_{\jepPubBold{Q}} \jepPubBold{C}$ and ${\rm DR}(\jepPubScript{M})$, i.e. 
the perverse sheaf corresponding to $\jepPubScript{M}$ via the Riemann-Hilbert correspondence.
\item[(iv)] $W$ is a finite, increasing filtration on $(\jepPubScript{M},F, \jepPubScript{P},\varphi)$, the \emph{weight filtration} of the mixed Hodge module.
\end{enumerate}
For a such an object to be a mixed Hodge module, it has to satisfy 
a complicated set of conditions of an inductive nature, which we do not discuss here.
The main reference for the basic definitions and results of this theory is \cite{Saito-MHM}; see also 
\cite{Saito-YPG} for an introduction. 


Given a mixed Hodge module $(\jepPubScript{M},F,\jepPubScript{P},\varphi, W)$, we say that the filtered $\jepPubScript{D}$-module $(\jepPubScript{M},F)$ is a \emph{Hodge $\jepPubScript{D}$-module}
(or that it
\emph{underlies} a mixed Hodge module). In fact, this is the only piece of information that we will be concerned with in this article. 
The basic example of a mixed Hodge module is ${\mathbf Q}_X^H[n]$, the trivial one. In this case, the filtered $\jepPubScript{D}$-module is the structure sheaf $\jepPubScript{O}_X$, 
with the filtration such that ${\rm gr}_p^F\jepPubScript{O}_X=0$ for all $p\neq 0$. The corresponding perverse sheaf is $\jepPubBold{Q}_X [n]$ and the weight filtration is such that 
${\rm gr}_p^W\jepPubScript{O}_X=0$ for $p\neq n$. 

The mixed Hodge modules on $X$ form an Abelian category, denoted ${\rm MHM}(X)$. Morphisms in this category are strict with respect to both the Hodge
and the weight filtration. The corresponding bounded derived category is denoted ${\bf D}^b\big({\rm MHM}(X)\big)$. 

Mixed Hodge modules satisfy Grothendieck's 6 operations formalism. The relevant fact for us is that to every morphism $f\colon X\to Y$ of 
smooth complex algebraic varieties we have a corresponding push-forward functor $f_+\colon {\bf D}^b\big({\rm MHM}(X)\big)\rightarrow {\bf D}^b\big({\rm MHM}(Y)\big)$
(this is denoted by $f_*$ in \cite{Saito-MHM}). Moreover, if
$g\colon Y\to Z$ is another such morphism, we have a functorial isomorphism $(g\circ f)_+\simeq g_+\circ f_+$. 

Regarding the push-forward functor for mixed Hodge modules, we note that on the level of $\jepPubScript{D}$-modules, it coincides with the usual $\jepPubScript{D}$-module push-forward. 
Moreover, if $f\colon X\to Y$ is proper and if we denote by ${\rm FM}(\jepPubScript{D}_X)$ the category of filtered $\jepPubScript{D}$-modules on $X$ (here it is convenient to work with right $\jepPubScript{D}$-modules), then
Saito defined in \cite{Saito-MHP} a functor
$$f_+ \colon {\bf D}^b \big({\rm FM}(\jepPubScript{D}_X)\big) \longrightarrow {\bf D}^b \big({\rm FM}(\jepPubScript{D}_Y)\big).$$
This is compatible with the usual direct image functor for right $\jepPubScript{D}$-modules and it is used to define
the push-forward between the derived categories of mixed Hodge modules at the level of filtered complexes. 
With a slight abuse of notation, if $(\jepPubScript{M},F)$ underlies a mixed Hodge module $M$ on $X$ and if $f\colon X\to Y$
is an arbitrary morphism, then we write $f_+(\jepPubScript{M},F)$ for the object in ${\bf D}^b\big({\rm FM}(\jepPubScript{D}_Y)\big)$ underlying $f_+M$.


An important feature of the push-forward of Hodge $\jepPubScript{D}$-modules with respect to proper morphisms is \emph{strictness}. 
This says that if $f\colon X\to Y$ is proper and $(\jepPubScript{M},F)$ underlies a mixed Hodge module on $X$, then 
$f_+ (\jepPubScript{M}, F)$ is strict 
as an object in ${\bf D}^b \big({\rm FM}(\jepPubScript{D}_Y)\big)$ (and moreover, each $H^i f_+ (\jepPubScript{M}, F)$ underlies a Hodge $\jepPubScript{D}_Y$-module). 
This means that the natural mapping 
\begin{equation}\label{strictness_formula}
R^i f_* \big(F_k (\jepPubScript{M} \overset{\jepPubBold{L}}{\otimes}_{\jepPubScript{D}_X} \jepPubScript{D}_{X\to Y}) \big) \longrightarrow 
R^i f_* (\jepPubScript{M} \overset{\jepPubBold{L}}{\otimes}_{\jepPubScript{D}_X} \jepPubScript{D}_{X\to Y})
\end{equation}
is injective for every $i, k\in\jepPubBold{Z}$. 
Taking $F_kH^if_+(\jepPubScript{M},F)$ to be the image of this map, we get the filtration on $H^if_+(\jepPubScript{M},F)$.

The push-forward with respect to open embeddings is more subtle. For example, suppose that $Z$ is an effective divisor
on the smooth variety $X$ and $j\colon U=X\smallsetminus Z\hookrightarrow X$ is the corresponding open immersion. 
Recall that $\jepPubScript{O}_X(*Z)$
is the push-forward $j_*\jepPubScript{O}_U$; on a suitable affine open neighborhood $V$ of a given point in $X$, this is given by localizing $\jepPubScript{O}_X(V)$ at 
an equation defining $Z\cap V$ in $V$.
$\jepPubScript{O}_X(*Z)$ has a natural left $\jepPubScript{D}$-module structure induced by the canonical $\jepPubScript{D}$-module structure on $\jepPubScript{O}_X$. In fact, as such
we have $\jepPubScript{O}_X(*Z)\simeq j_+\jepPubScript{O}_U$  (in general, for a $\jepPubScript{D}_U$-module $\jepPubScript{M}$, the $\jepPubScript{D}$-module push-forward
 $j_+\jepPubScript{M}$ agrees with $j_*\jepPubScript{M}$, with the induced $\jepPubScript{D}_X$-module structure).
We thus see that $\jepPubScript{O}_X(*Z)$ carries a canonical filtration such that the corresponding
filtered $\jepPubScript{D}$-module underlies $j_+\jepPubBold{Q}_U^H[n]$. This filtration is the one that leads to the Hodge ideals studied in \cite{MP1}.



